\chapter{Memory}
\label{app:memory}

\section{How much there is}
\index{memory}

The Commander X16 ships with 512K of banked memory and can be expanded to 2M.
\ccprod{} uses all of it for the workbook, and the figure on the status line at
startup says how much that is on your machine.

\begin{cctable}{Machine}{Free at startup}{2.4in}{2.42in}
Standard, 512K & about 516,000 bytes \\
Fully expanded, 2M & about 2,097,000 bytes \\
\end{cctable}

That figure is a statement of capacity, not of use. It is measured once and does
not fall as the workbook grows.

\section{What a workbook costs}
\index{memory!what a workbook costs}

\begin{cctable}{Item}{Costs}{2.4in}{2.42in}
One cell, of any kind & 8 bytes \\
One worksheet, whether used or not & 1K \\
Each block of 256 rows you put anything in & 2K \\
A piece of text & its length, plus a few bytes \\
A formula & its text, and its compiled form \\
\end{cctable}

The third row is the one that catches people out. Storage is allocated in bands
of 256 rows, and a band costs the same whether it holds one cell or two hundred
and fifty-six.

\begin{ccnote}
A worksheet using rows 1 to 100 costs about 3K. A worksheet with a single cell
scattered into each of 256 different bands costs 513K --- the same data, spread
out, for a hundred and seventy times the memory.

If a workbook will not fit, look first at how far apart its rows are.
\end{ccnote}

At the program's own limits --- 16,384 cells on each of sixteen worksheets ---
the cells alone come to 128K and the worksheet indexes to another 16K, which is
why those limits are reachable on an expanded machine and not on a standard one.

\section{When memory runs out}

\begin{cctable}{Message}{Means}{1.8in}{3.02in}
\msg{Out of memory} & There is no more banked memory. A larger workbook needs a
machine with more of it. \\
\msg{Resource limit reached} & Memory is not the problem: a fixed limit has been
reached --- 16,384 cells on a worksheet, sixteen worksheets, 128 number formats,
8,192 pieces of text. Appendix~\ref{app:limits} lists them all. \\
\end{cctable}

An import that runs out does not fail. It stops, keeps what it has, and reports
that not everything fitted.

\section{The address space}
\index{memory!map}

For completeness, and because it explains why a program this small needs banked
memory at all: of the 64K the processor can address, a program on the Commander
X16 gets about 30K and the rest belongs to the machine.

\begin{ccfigure}
\begin{tikzpicture}[x=1pt,y=-1pt,
  reg/.style={draw=ccink,line width=0.6pt,minimum width=118pt,inner sep=2pt,
              font=\sffamily\fontsize{7.6}{9.2}\selectfont,anchor=north west},
  adr/.style={font=\ttfamily\fontsize{7.4}{9}\selectfont,text=ccink,
              anchor=north east},
  ann/.style={font=\sffamily\fontsize{7.2}{8.8}\selectfont,text=ccink,
              anchor=west}]

  \node[reg,minimum height=15pt,fill=black!22] at (46,0)
    {\$0000--\$07FF\ \ the machine's own workspace};
  \node[reg,minimum height=60pt,fill=black!6] at (46,15)
    {\$0801--\$7FFF\ \ \textbf{the program}};
  \node[reg,minimum height=30pt,fill=black!14] at (46,75)
    {\$8000--\$9EFF\ \ \textbf{loaded as needed}};
  \node[reg,minimum height=15pt,fill=black!30] at (46,105)
    {\$9F00--\$9FFF\ \ the display and other hardware};
  \node[reg,minimum height=32pt,fill=black!14] at (46,120)
    {\$A000--\$BFFF\ \ \textbf{a window on banked memory}};
  \node[reg,minimum height=30pt,fill=black!22] at (46,152)
    {\$C000--\$FFFF\ \ a window on the machine's ROM};

  \node[adr] at (42,-1)  {\$0000};
  \node[adr] at (42,70)  {\$7FFF};
  \node[adr] at (42,100) {\$9EFF};
  \node[adr] at (42,147) {\$BFFF};
  \node[adr] at (42,177) {\$FFFF};

  \node[ann] at (170,22) {The display, the keyboard, the cells,};
  \node[ann] at (170,33) {the workbook. Always present.};

  \node[ann] at (170,82) {The parts of the program that are};
  \node[ann] at (170,93) {read from the card when wanted.};

  \node[ann] at (170,127) {8K at a time out of 512K to 2M.};
  \node[ann] at (170,138) {The whole workbook lives here.};

  \node[ann] at (170,160) {All arithmetic is done by the};
  \node[ann] at (170,171) {machine's own maths routines.};
\end{tikzpicture}
\caption{The address space as \ccprod{} uses it}
\label{fig:memmap}
\end{ccfigure}

The workbook does not live in that 30K --- it lives in banked memory, reached 8K
at a time through the window at \$A000. The program does not fit in it either,
which is why the block at \$8000 is filled from the card with whichever part of
the program is wanted next.

\begin{ccnote}
That is the reason for two things you will notice in ordinary use. Commands
pause before they start, because the machine is reading from the card. And the
card has to stay in the drive for the whole session, not just at startup.
\end{ccnote}

\section{If the status line says STACK!}
\index{STACK!}

\begin{cccaution}
The program has run out of its own internal working space. It noticed rather
than crashing, but nothing after that point is guaranteed.

Save the workbook \textbf{under a new name} --- not over your existing file ---
and restart \ccprod{}. Then check the saved workbook before trusting it.
\end{cccaution}
